\ZSFNewColumn
\StartChapter{Integrale\ZSFScriptRef{70}}

\SubsectionBar{Integrationsregeln\ZSFScriptRef{71}}
\ZSFindex{Integrationsregeln}
\ZSFindex{Hauptsatz der Integralrechnung}
\begin{tablebox}
\begin{ZSFtable}[header=false]{Y{0.7} Z{1.3}}
    \ZSFlabel{Linearität} &
    $\displaystyle \int_a^b (f(x)+g(x))\,dx = \int_a^b f(x)\,dx + \int_a^b g(x)\,dx$ \\
    \ZSFlabel{Ableitung mit variablen Grenzen} &
    $\displaystyle \frac{d}{dx}\!\left(\int_{a(x)}^{b(x)} f(t)\,dt\right)=b'(x)f(b(x))-a'(x)f(a(x))$ \\
    \ZSFlabel{Additivität der Intervalle} &
    $\displaystyle \int_a^b f(x)\,dx + \int_b^c f(x)\,dx = \int_a^c f(x)\,dx$ \\
\end{ZSFtable}
\end{tablebox}

\begin{tablebox}
\begin{ZSFtable}{C C}
    \ZSFheaderRow \ZSFhead{Gerade Funktionen} & \ZSFhead{Ungerade Funktionen} \\
    $\displaystyle \int_{-a}^{0} f(x)\,dx = \int_0^a f(x)\,dx$ &
    $\displaystyle \int_{-a}^{0} f(x)\,dx = -\int_0^a f(x)\,dx$ \\
\end{ZSFtable}
\end{tablebox}

\SubsectionBar{Uneigentliche Integrale\ZSFScriptRef{72}}
Mit Hilfe der \ZSFkeyword[uneigentliches Integral@uneigentliches Integral]{uneigentlichen Integrale} ist es möglich, Funktionen zu integrieren deren Definitionsbereich unbeschränkt ist:
\textVorBox{1. Ordnung (Polstelle / Definitionslücke bei $a$):}
\begin{formulabox}
    $\displaystyle \int_a^b f(x)\,dx = \ZSFlimAuto{\zeta \to a^+}\int_{\zeta}^{b} f(x)\,dx$
\end{formulabox}
\textNachBox{$\zeta$ ist die bewegliche Grenze und nähert sich der Problemstelle $a$.}
\textVorBox{2. Ordnung (unbeschränkter Definitionsbereich):}
\begin{formulabox}
    $\displaystyle \int_a^{\infty} f(x)\,dx = \ZSFlimAuto{\zeta \to \infty}\int_a^{\zeta} f(x)\,dx$
\end{formulabox}

\SubsectionBar{Hilfsmittel\ZSFScriptRef{71}}
\textVorBox{\ZSFkeyword{Partielle Integration}:}
\begin{formulabox}
    $\displaystyle \int_a^b u'(x)v(x)\,dx = [u(x)v(x)]_a^b - \int_a^b u(x)v'(x)\,dx$
\end{formulabox}

Tricks, um Integrale mit partieller Integration zu lösen:
\begin{ZSFlist}
    \ZSFItem{Term 2-mal partiell integrieren.}
    \ZSFItem{Setze $u'(x)=1$ (respektive $v'(x)=1$).}
\end{ZSFlist}

\textVorBox{\ZSFkeyword{Partialbruchzerlegung}:}
\begin{tablebox}
\begin{ZSFtable}[header=false]{Y{1} C}
    Einfache Nullstelle $x_0$: & $\displaystyle \frac{A}{x-x_0}$ \\
    Zweifache Nullstelle $x_0$: & $\displaystyle \frac{A}{(x-x_0)^2} + \frac{B}{x-x_0}$ \\
    Komplexe Nullstelle: & $\displaystyle \frac{Ax+B}{x^2+px+q}$ \\
\end{ZSFtable}
\end{tablebox}

\textVorBox{Beispiel:}
\begin{formulabox}
    \[\frac{x}{x^3+x^2-x-1} = \frac{A}{(x+1)^2} + \frac{B}{x+1} + \frac{C}{x-1}\]
    $x=0\colon\; {-A-B+C=0}$\quad
    $x=1\colon\; {4C=1}$\quad
    $x=2\colon\; {A+3B+9C=2}$\\[2pt]
    $\Rightarrow\; A=\dfrac{1}{2},\quad B=-\dfrac{1}{4},\quad C=\dfrac{1}{4}$
\end{formulabox}
\textNachBox{\ZSFdanger{Tipp:} $(1+x^3)=(1+x)(1-x+x^2)$}

\ZSFNewColumn
\SubsectionBar{Substitution\ZSFScriptRef{71}}
\ZSFindex{Substitution}
\textVorBox{Bei der \ZSFkeyword{Substitution} ersetzt man komplizierte Terme. Im folgenden Beispiel wird $f(x)$ durch $u$ ersetzt (Heuristik \& Standard-Ansätze siehe Anhang \ZSFref{sec:substitutions_heuristik}):}
\begin{formulabox}
    $\displaystyle \int_{\ZSFmhlA{a}}^{\ZSFmhlA{b}} f'(x)g'(f(x))\,\ZSFmhlA{dx} = \int_{\ZSFmhlB{f(a)}}^{\ZSFmhlB{f(b)}} g'(u)\,\ZSFmhlB{du}$
\end{formulabox}

\SubsectionBar{Anwendungen\ZSFScriptRef{75,76}}
\SubsubsectionBar{Flächen, Sektorfläche und Bogenlänge}
\ZSFindex{Fläche unter einer Kurve}
\textVorBox{Das Integral liefert die orientierte Fläche zwischen einer Kurve und der $x$-Achse:}
\begin{formulabox}
    \ZSFlabel{Explizit:}\quad $\displaystyle A_{\mathrm{or}} = \int_a^b f(x)\,dx$\par
    \ZSFsep
    \ZSFlabel{Parametrisiert:}\quad $\displaystyle A_{\mathrm{or}} = \int_{t_1}^{t_2} \dot x(t)y(t)\,dt$\par
\end{formulabox}
\textNachBox{Für die geometrische Fläche gilt $A=\int_a^b|f(x)|\,dx$. Liegt der Graph vollständig oberhalb der Achse, kann der Betrag entfallen; bei Achsenkreuzungen in Teilintervalle zerlegen. Parametrisiert gilt dieselbe Regel: Orientierung beachten und an Achsenkreuzungen trennen.}

\textVorBox{Die orientierte Fläche zwischen einer Kurve und der $y$-Achse beträgt:}
\begin{formulabox}
    $\displaystyle \text{Parametrisiert:}\quad A_{\mathrm{or},y} = \int_{t_1}^{t_2} x(t)\dot y(t)\,dt$
\end{formulabox}
\textNachBox{Für die geometrische Fläche Orientierung beachten und bei Kreuzungen der $y$-Achse in Teilflächen zerlegen.}

\textVorBox{Die \ZSFkeyword{Sektorfläche} ist die Fläche zwischen dem Ursprung und zwei Punkten einer Kurve:}
\begin{formulabox}
    $\displaystyle A=\pm\frac{1}{2}\int_{t_1}^{t_2}\!\left(x(t)\dot y(t)-\dot x(t)y(t)\right)\,dt
    =\pm\frac{1}{2}\int_{\varphi_1}^{\varphi_2}\!\rho(\varphi)^2\,d\varphi$
\end{formulabox}
\textNachBox{positiv, wenn die Fläche links der Kurve ist und negativ, wenn sie rechts der Kurve ist.}

\textVorBox{Die Bogenlänge einer Kurve oder Funktion:}
\begin{formulabox}
    \ZSFlabel{Explizit:}\quad $\displaystyle s = \int_a^b \sqrt{1+f'(x)^2}\,dx$\par
    \ZSFsep
    \ZSFlabel{Parametrisiert:}\quad $\displaystyle s = \int_{t_1}^{t_2}\sqrt{\dot x(t)^2+\dot y(t)^2}\,dt$\par
    \ZSFsep
    \ZSFlabel{Polarkoordinaten:}\quad $\displaystyle s = \int_{\varphi_1}^{\varphi_2}\sqrt{\rho(\varphi)^2+\dot \rho(\varphi)^2}\,d\varphi$\par
\end{formulabox}

\SubsubsectionBar{Rotationsvolumen und Rotationsoberfläche}
\textVorBox{Das \ZSFkeyword{Rotationsvolumen} um die $x$-Achse beträgt:}
\begin{formulabox}
    \ZSFlabel{Explizit:}\quad $\displaystyle V_x = \pi\int_a^b f(x)^2\,dx$\par
    \ZSFsep
    \ZSFlabel{Parametrisiert:}\quad $\displaystyle V_x = \pi\left|\int_{t_1}^{t_2}\dot x(t)y^2(t)\,dt\right|$\par
\end{formulabox}
\textNachBox{wobei $a$ und $b$ zwei Punkte auf der $x$-Achse liegen.}

\textVorBox{Das \ZSFkeyword{Rotationsvolumen} um die $y$-Achse beträgt:}
\begin{formulabox}
    \ZSFlabel{Explizit (Fläche zwischen Graphen und y-Achse):}\quad $\displaystyle V_y = \pi\int_c^d f(y)^2\,dy$\par
    \ZSFsep
    \ZSFlabel{Umgeschrieben:}\quad $\displaystyle V_y = \pi\int_a^b x^2\cdot f'(x)\,dx$\par
    \ZSFsep
    \ZSFlabel{Explizit (Fläche zwischen Graphen und x-Achse):}\quad $\displaystyle V_y = 2\pi\int_a^b xf(x)\,dx$\par
    \ZSFsep
    \ZSFlabel{Parametrisiert:}\quad $\displaystyle V_y = \pi\left|\int_{t_1}^{t_2}x^2(t)\dot y(t)\,dt\right|$\par
\end{formulabox}
\textNachBox{$a,b$: x-Grenzen; $c,d$: y-Grenzen; $t_1,t_2$: Parametergrenzen.}

\textVorBox{Die \ZSFkeyword{Rotationsoberfläche} um die $x$-Achse beträgt:}
\begin{formulabox}
    \ZSFlabel{Explizit:}\quad $\displaystyle O_x = 2\pi\int_a^b f(x)\sqrt{1+f'(x)^2}\,dx$\par
    \ZSFsep
    \ZSFlabel{Parametrisiert:}\quad $\displaystyle O_x = 2\pi\int_{t_1}^{t_2}|y(t)|\sqrt{\dot x(t)^2+\dot y(t)^2}\,dt$\par
\end{formulabox}

\SubsubsectionBar{Schwerpunkt und Trägheitsmomente}
\textVorBox{Der \ZSFkeyword{Schwerpunkt} $(x_S,y_S)$ ist der Kräftemittelpunkt:}
\begin{formulabox}
    \ZSFlabel{x-Koordinate:}\quad $\displaystyle x_S\int_a^b G(x)\,dx = x_SA = \int_a^b xG(x)\,dx$\par
    \ZSFsep
    \ZSFlabel{y-Koordinate:}\quad $\displaystyle y_S\int_c^d H(y)\,dy = y_SA = \int_c^d yH(y)\,dy$\par
\end{formulabox}
\textNachBox{wobei $G(x)$ und $H(y)$ die Ausdehnung in y- und x-Richtung sind, $a$ und $b$ auf der x-Achse und $c$ und $d$ auf der y-Achse liegen.}

\textVorBox{Das \ZSFkeyword{Massenträgheitsmoment} gibt die Trägheit eines starren Körpers gegenüber einer Änderung seiner Winkelgeschwindigkeit während der Drehung um eine bestimmte Achse an:}
\begin{formulabox}
    \ZSFformulaline{\theta_y = \int_a^b \rho(x)x^2G(x)\,dx}{Massenträgh.}
\end{formulabox}

\textNachBox{wobei $\rho$ die Dichte ist.}
\ZSFindexsee{Trägheitsmoment}{Massenträgheitsmoment}
\textVorBox{Für $\rho=1$ sei das \ZSFkeyword{Flächenträgheitsmoment}:}
\begin{formulabox}
    \ZSFlabel{x-Achse:}\quad $\displaystyle I_x = \frac{1}{3}\int_a^b f(x)^3\,dx = \int_c^d y^2H(y)\,dy$\par
    \ZSFsep
    \ZSFlabel{y-Achse:}\quad $\displaystyle I_y = \int_a^b x^2f(x)\,dx$\par
\end{formulabox}
\textNachBox{$a,b$: x-Grenzen; $c,d$: y-Grenzen.}

\textVorBox{Das polare Flächenträgheitsmoment einer Kreisscheibe beträgt:}
\begin{formulabox}
    \ZSFformulaline{I_0 = \frac{1}{2}\pi r^4}{Kreisscheibe}
\end{formulabox}

\textVorBox{Das Trägheitsmoment eines rotierenden Graphen um die $x$-Achse:}
\begin{formulabox}
    $\displaystyle \theta_x = \frac{\pi\rho}{2}\int_a^b f(x)^4\,dx$
\end{formulabox}
\textNachBox{Die Berechnung für die $y$- und $z$-Achse erfolgt analog.}

\textVorBox{Die \ZSFkeyword{kinetische Energie} der Rotation um eine Achse:}
\begin{formulabox}
    \ZSFformulaline{T = \frac{1}{2}\theta_{x,y,z}\omega^2}{kinet. Energie}
\end{formulabox}
